% !TEX program = lualatex
\documentclass[lualatex,article,paper=a4paper,jafontsize=10pt]{jlreq}
\usepackage[top=15mm,bottom=15mm,left=12mm,right=12mm]{geometry}
\usepackage{luatexja-fontspec}
\usepackage{multirow}
\usepackage{array}

% 日本語フォント（Mac の Hiragino Kaku Gothic ProN）
\setmainjfont{Hiragino Kaku Gothic ProN W3}[BoldFont=Hiragino Kaku Gothic ProN W6]
\setsansjfont{Hiragino Kaku Gothic ProN W6}
% 欧文フォントも和文と統一（数字 111/121 等を表中で違和感なく）
\setmainfont{Hiragino Kaku Gothic ProN W3}[BoldFont=Hiragino Kaku Gothic ProN W6]
\setsansfont{Hiragino Kaku Gothic ProN W6}

\pagestyle{empty}
\renewcommand{\arraystretch}{1.2}
\setlength{\tabcolsep}{4pt}

\begin{document}

\begin{center}
{\Large\textbf{表1\quad WBS（作業分担表）}}\\[0.3em]
{\small ハッカソン1 チーム23 — 議事録 表1 の再設計版（v0.1, 2026-04-27 塩澤 匠生）}
\end{center}

\vspace{0.3em}

%==================================================
% 100 設計フェーズ
%==================================================
\noindent
\begin{tabular}{|c|c|l|p{9cm}|c|}
\hline
\textbf{時間的フェーズ} & \textbf{番号} & \textbf{要素成果物} & \textbf{活動タスク} & \textbf{担当者} \\
\hline
\multirow{12}{*}{\shortstack{100\\設計\\フェーズ}}
 & \multirow{7}{*}{110} & \multirow{7}{*}{基本設計}
   & 111 FBS（機能分解構造）の最終化 & 全員 \\\cline{4-5}
 & & & 112 PBS（製品分解構造）の最終化 & 全員 \\\cline{4-5}
 & & & 113 MOE / MOP / TPM の最終化 & 全員 \\\cline{4-5}
 & & & 114 システムアーキテクチャ設計（EMA準拠） & 塩澤 \\\cline{4-5}
 & & & 115 通信プロトコル基本設計（CTRL/BEAT/NOTE） & 塩澤 \\\cline{4-5}
 & & & 116 楽譜データ形式 基本設計（4パート/ROM埋込） & 齋藤 \\\cline{4-5}
 & & & 117 音色合成 基本設計（金管倍音 + ADSR） & 梅澤 \\\cline{2-5}
 & \multirow{5}{*}{120} & \multirow{5}{*}{詳細設計}
   & 121 共通層API詳細（IModule / SystemData / ProjectConfig） & 塩澤 \\\cline{4-5}
 & & & 122 node\_01 詳細設計（IMU→拍検出→テンポ推定→送出） & 塩澤 \\\cline{4-5}
 & & & 123 node\_02-05 詳細設計（受信→楽譜進行→NOTE送出） & 塩澤 \\\cline{4-5}
 & & & 124 Processing 詳細設計（NOTE受信→ボイス管理→出力） & 梅澤 \\\cline{4-5}
 & & & 125 楽譜データ詳細設計（課題曲のパート割・データ表化） & 齋藤 \\
\hline
\end{tabular}

\vspace{0.8em}

%==================================================
% 200 製造フェーズ
%==================================================
\noindent
\begin{tabular}{|c|c|l|p{9cm}|c|}
\hline
\textbf{時間的フェーズ} & \textbf{番号} & \textbf{要素成果物} & \textbf{活動タスク} & \textbf{担当者} \\
\hline
\multirow{19}{*}{\shortstack{200\\製造\\フェーズ}}
 & \multirow{3}{*}{210} & \multirow{3}{*}{共通層実装}
   & 211 IModule / ModuleTimer 実装 & 塩澤 \\\cline{4-5}
 & & & 212 SystemData / ProjectConfig テンプレ整備 & 塩澤 \\\cline{4-5}
 & & & 213 OrcNetModule（UDP送受信ラッパ：CTRL/BEAT/NOTE 共通） & 塩澤 \\\cline{2-5}
 & \multirow{5}{*}{220} & \multirow{5}{*}{\shortstack[l]{指揮者ノード\\(node\_01)}}
   & 221 IMUドライバモジュール & 塩澤 \\\cline{4-5}
 & & & 222 信号前処理（重力分離・LPF）モジュール & 塩澤 \\\cline{4-5}
 & & & 223 拍検出モジュール（候補方式比較含む） & 塩澤 \\\cline{4-5}
 & & & 224 テンポ推定モジュール & 塩澤 \\\cline{4-5}
 & & & 225 コマンド送出モジュール（CTRL/BEAT） & 塩澤 \\\cline{2-5}
 & \multirow{4}{*}{230} & \multirow{4}{*}{\shortstack[l]{楽器ノード\\(node\_02-05)}}
   & 231 CTRL/BEAT 受信モジュール & 塩澤 \\\cline{4-5}
 & & & 232 楽譜データ保持・進行ロジック & 塩澤 \\\cline{4-5}
 & & & 233 NOTE送出モジュール & 塩澤 \\\cline{4-5}
 & & & 234 パート別差分適用（node\_03/04/05 への展開） & 塩澤 \\\cline{2-5}
 & \multirow{3}{*}{240} & \multirow{3}{*}{楽譜データ準備}
   & 241 課題曲 4 パートの選定 & 齋藤 \\\cline{4-5}
 & & & 242 4 パート楽譜のヘッダ配列化（PROGMEM 形式） & 齋藤 \\\cline{4-5}
 & & & 243 楽譜データの ROM 埋込確認 & 齋藤 \\\cline{2-5}
 & \multirow{4}{*}{250} & \multirow{4}{*}{Processing 側実装}
   & 251 NOTE 受信モジュール & 梅澤 \\\cline{4-5}
 & & & 252 金管音色合成エンジン（倍音 + ADSR） & 梅澤 \\\cline{4-5}
 & & & 253 4 パート同時発音管理（ボイスマネージャ） & 梅澤 \\\cline{4-5}
 & & & 254 UI / モード表示（最小限） & 梅澤 \\
\hline
\end{tabular}

\vspace{0.8em}

%==================================================
% 300 テストフェーズ
%==================================================
\noindent
\begin{tabular}{|c|c|l|p{9cm}|c|}
\hline
\textbf{時間的フェーズ} & \textbf{番号} & \textbf{要素成果物} & \textbf{活動タスク} & \textbf{担当者} \\
\hline
\multirow{15}{*}{\shortstack{300\\テスト\\フェーズ}}
 & \multirow{4}{*}{310} & \multirow{4}{*}{単体テスト}
   & 311 共通層単体（OrcProtocol / ModuleTimer） & 塩澤 \\\cline{4-5}
 & & & 312 拍検出ゴールデンテスト（記録 IMU 波形→期待発火） & 塩澤 \\\cline{4-5}
 & & & 313 楽譜進行 状態機械テスト & 塩澤 \\\cline{4-5}
 & & & 314 音色生成 単体確認（聴感+波形目視） & 梅澤 \\\cline{2-5}
 & \multirow{4}{*}{320} & \multirow{4}{*}{結合テスト}
   & 321 node\_01 単体デモ（手振り→ログ拍出力） & 塩澤 \\\cline{4-5}
 & & & 322 node\_01 + node\_02 結合（BEAT 送受信） & 塩澤 \\\cline{4-5}
 & & & 323 4 ノード同時 同期誤差計測 & 全員 \\\cline{4-5}
 & & & 324 全ノード + Processing 全結合（実音出し） & 全員 \\\cline{2-5}
 & \multirow{7}{*}{330} & \multirow{7}{*}{システムテスト}
   & 331 TP-1 拍検出精度（80 / 120 / 160 BPM） & 塩澤 \\\cline{4-5}
 & & & 332 TP-2 通信遅延 & 塩澤 \\\cline{4-5}
 & & & 333 TP-3 同期誤差（4 ノード時刻合わせ） & 全員 \\\cline{4-5}
 & & & 334 TP-4 テンポ追従（80→120 BPM 階段変化） & 塩澤 \\\cline{4-5}
 & & & 335 TP-5 パケロス耐性（5\% 擬似パケロス） & 塩澤 \\\cline{4-5}
 & & & 336 TP-6 CPU 負荷（入力フェーズ 2ms 以下） & 塩澤 \\\cline{4-5}
 & & & 337 TP-7 起動時間（電源→演奏可能 5s 以下） & 塩澤 \\
\hline
\end{tabular}

\vspace{0.8em}

%==================================================
% 400 ストレッチフェーズ
%==================================================
\noindent
\begin{tabular}{|c|c|l|p{9cm}|c|}
\hline
\textbf{時間的フェーズ} & \textbf{番号} & \textbf{要素成果物} & \textbf{活動タスク} & \textbf{担当者} \\
\hline
\multirow{3}{*}{\shortstack{400\\ストレッチ\\フェーズ}}
 & \multirow{3}{*}{410} & \multirow{3}{*}{強弱推定}
   & 411 振幅 → velocity 推定アルゴ実装 & 塩澤 \\\cline{4-5}
 & & & 412 CTRL パケットへの velocity 乗せ込み & 塩澤 \\\cline{4-5}
 & & & 413 TP-8 強弱制御テスト（3 段階分離） & 塩澤 \\
\hline
\end{tabular}

\vspace{0.8em}

%==================================================
% 500 運用・発表
%==================================================
\noindent
\begin{tabular}{|c|c|l|p{9cm}|c|}
\hline
\textbf{時間的フェーズ} & \textbf{番号} & \textbf{要素成果物} & \textbf{活動タスク} & \textbf{担当者} \\
\hline
\multirow{2}{*}{\shortstack{500\\運用\\発表}}
 & 510 & 発表 & 511 発表準備・デモ調整 & 全員 \\\cline{2-5}
 & 520 & 報告書 & 521 成果報告書（個人） & 全員 \\
\hline
\end{tabular}

\vspace{1em}

\noindent
\textbf{担当者表記}：「塩澤」=25G1065 塩澤匠生、「齋藤」=25G1053 齋藤翔太、「梅澤」=25G1021 梅澤颯太、「全員」=実装3人全員（塩澤・齋藤・梅澤）。\\
24G勢（片岡・地曵・御代川）は議事録持ち回りと 300 テストフェーズ以降の実機検証・発表準備に合流する想定のため、本表の活動タスク列には記載しない。

\end{document}
